\section{The degree-\(163\) arithmetic image}
\label{sec:arithmetic-cm-163}

Put \(p=163\), \(s=[2]\), \(h=81\), and \(c=s^{81}\).
The residue \(2\) has order \(162\) modulo \(163\);
therefore
\(\mathcal O=\mathbb Z_2[\zeta_{163}]\) is unramified of
degree \(162\). The group
\[
 T=\{u\in1+2\mathcal O:u\,c(u)=1\}
\]
has rank \(81\) and torsion \(\{\pm1\}\).
Since \(s-1\) kills this torsion and is invertible on the
rational anti-invariant space,
\[
 P=(s-1)T\simeq\mathbb Z_2^{81}.
\]
For a reproducible description of the reflex polynomial,
write \(Q=\sum_{j\in I_+}s^j-\sum_{j\in I_-}s^j\), where
\[
\begin{split}
 I_+={}&\{2,4,9,13,16,18,20,25,27,31,35,42,44,49,
                                     52,54,57,67,70,73\},\\
 I_-={}&\{1,3,5,11,15,17,19,22,26,28,32,38,43,47,50,
                                  53,56,62,68,71,75\}.
\end{split}
\]
Then \(Q(1)=-1\), and exact elimination on the signed
cyclic lattice gives
\[
 \det Q=40176902897933369663096255063282475371772362696316,
 \qquad v_2(\det Q)=2.
\]
The corresponding mod-two calculation is
\[
 \gcd(\overline Q,s^{81}+1)=s^2+s+1.
\]
Because \(Q\) is an automorphism on the only sign factor,
\begin{equation}
 [T:QT]=[P:QP]=4,\qquad
 T/QT\xrightarrow{\ \theta\ }
      (\mathcal O/2\mathcal O)^c/\overline Q
                       (\mathcal O/2\mathcal O)^c
 \quad\text{is an isomorphism}.
 \label{eq:163-first-layer}
\end{equation}
The last assertion follows from surjectivity of the
first-layer map and equality of the two orders.
Thus
\[
 \mathscr C=P/QP\simeq
     \mathbb F_2[s]/(s^2+s+1),
\]
with the irreducible cyclotomic \(C_3\)-action.

\begin{lemma}
\label{lem:163-period-unit-square}
Let \(H_0=\langle8\rangle\subset(\mathbb Z/163)^\times\),
and put
\[
 \theta_i=\sum_{a\in2^iH_0}\zeta_{163}^a,\qquad i=0,1,2.
\]
Then
\begin{equation}
 \prod_{j=0}^{26}
      (\zeta_{163}^{8^j}+\zeta_{163}^{-8^j})
            =-(5+\theta_1)^2.
 \label{eq:163-period-square}
\end{equation}
The element \(5+\theta_1\) is a real unit.
The two conjugate identities hold for the other two cosets.
\end{lemma}
\begin{proof}
The subgroup \(H_0\) has order \(54\), contains \(-1\),
and its three cosets partition the nonzero residues.
Thus the periods are real and
\(\theta_0+\theta_1+\theta_2=-1\).
Counting sums of two residues in these cosets gives
\begin{equation}
\begin{aligned}
 \theta_0^2&=54+20\theta_0+17\theta_1+16\theta_2,\\
 \theta_1^2&=54+16\theta_0+20\theta_1+17\theta_2,\\
 \theta_0\theta_1&=17\theta_0+16\theta_1+21\theta_2.
\end{aligned}
 \label{eq:163-period-table}
\end{equation}
Here every coefficient is an integer residue count.
For the product in \eqref{eq:163-period-square}, one can
work entirely in \(\mathbb Z[X]/(X^{163}-1)\), starting
with \(C_0(X)=1\) and using
\[
 C_{j+1}(X)=C_j(X)(X^{8^j}+X^{-8^j})
                     \pmod{X^{163}-1},\qquad0\leq j<27.
\]
The final coefficient is constant on each of
\(\{0\},H_0,2H_0,4H_0\), with respective values
\[
 823364,\qquad823427,\qquad823413,\qquad823426.
\]
These integers give
\[
 C_{27}(\zeta_{163})=-62+\theta_0-13\theta_1.
\]
The second line of \eqref{eq:163-period-table}, together
with the sum of the periods, gives
\[
 -(5+\theta_1)^2
  =-25-10\theta_1-\theta_1^2
  =-62+\theta_0-13\theta_1.
\]
This proves the identity. Every factor on its left is
a unit, since it is a root of unity times
\(1+\zeta_{163}^{2a}\) and \(\Phi_{163}(-1)=1\).
The algebraic integer \(5+\theta_1\) is consequently a
unit as well. Applying \(s\) and \(s^2\) proves the
conjugate identities.
\end{proof}

\begin{theorem}
\label{thm:163-exact-arithmetic-image}
After imposing any collection of the scalar kernels
\(h_q\) with \(q\ne163\), the degree-\(163\) arithmetic
images of the resulting \(R\)- and \(V\)-groups are
both exactly
\[
 Q_{163}P_{163},\qquad [P_{163}:Q_{163}P_{163}]=4.
\]
This is an unconditional integral image statement.
\end{theorem}
\begin{proof}
The exact first-layer hypothesis for
Theorem~\ref{thm:uniform-arithmetic-cm-image} holds by
\eqref{eq:163-first-layer}.
Write
\(\xi=s\bmod(s^2+s+1)\) in \(\mathscr C\).
For \(\ell\in\mathscr C^\vee\), the exponents
\(\ell(\xi^k)\) in the real-unit map are periodic of
period three. Put
\[
 U_i=\prod_{j=0}^{26}
  \bigl(\zeta_{163}^{2^{i+3j}}
                   +\zeta_{163}^{-2^{i+3j}}\bigr),
 \qquad i=0,1,2.
\]
The unit assigned to \(\ell\) is
\[
 u_\ell=U_0^{\ell(1)}U_1^{\ell(\xi)}U_2^{\ell(\xi^2)}.
\]
Lemma~\ref{lem:163-period-unit-square} and its conjugates
make every \(U_i\) a negative square of a real unit.
Thus \(\mathscr K_{163}\) is the zero map into real units
modulo signs and squares. Its kernel is all of
\(\mathscr C^\vee\), whose annihilator is zero.
The exact uniform arithmetic-image theorem now gives
the claimed image \(QP\). That theorem also provides
witnesses for every element of \(QP\) with all other
prime scalar coordinates equal to \(1\), so the equality
persists under any collection of the other-prime kernels.
\end{proof}

On a normalized pair \(R_0,V_0\) with pro-odd
\(R_0/V_0\), the full images of \(h_{163}\) coincide:
their quotient is both pro-odd and pro-\(2\), hence is
trivial. The common image contains \(QP\).
Its quotient in \(P/QP\) is stable under the full
cyclotomic action and this module is irreducible of
dimension two. Thus the common full image is either
\(QP\) or \(P\). The arithmetic equality alone does not
select either alternative.

\begin{lemma}
\label{lem:163-integral-pair-decomposition}
There is a cyclotomic-stable decomposition
\[
 P=B_2\oplus L_{79},\qquad QP=2B_2\oplus L_{79},
 \qquad
 \operatorname{rank}_{\mathbb Z_2}B_2=2,\quad
 \operatorname{rank}_{\mathbb Z_2}L_{79}=79.
\]
Moreover, if \(t=-s\) acts on the additive
\(\mathbb Z_2\)-lattice \(P\), then \(P\) is a free
rank-one \(\mathbb Z_2[C_{81}]\)-module.
\end{lemma}
\begin{proof}
Complex conjugation acts on the norm-one torus as
inversion, so \(s^{81}=-1\) on the additive lattice.
Consequently \(t=-s\) has order \(81\).
The rational representation is the signed cyclic
representation for \(s\), hence the regular cyclic
representation for \(t\), with characteristic
polynomial \(X^{81}-1\).

Since \(81\) is odd, \(\mathbb Z_2[C_{81}]\) is finite
etale over \(\mathbb Z_2\), a product of unramified
discrete valuation rings. The lattice \(P\) decomposes
under its integral idempotents. Each component is
torsion-free of rank one over its corresponding
valuation ring, because the rational representation
is regular. Each component is therefore free of
rank one, proving the module assertion.

Modulo \(2\), the only factor on which \(Q\) is
noninvertible is \(X^2+X+1\). It defines an unramified
degree-two component \(B_2\).
The exact determinant valuation is \(2\), so the
valuation of \(Q\) on this component is \(1\);
it is zero on every other component.
Therefore \(Q B_2=2B_2\), while \(Q\) is invertible
on the complementary rank-\(79\) lattice.
This proves the stated decomposition.
\end{proof}
